\providecommand{\BM}{\textbf{[B]}}
\providecommand{\KM}{\textbf{[K]}}
\providecommand{\SM}{\textbf{[S]}}
\providecommand{\XM}{\textbf{[X]}}
\providecommand{\QM}{\textbf{[Q]}}

\chapter*{Der Koeffizient 8$\pi$ und der Status von Issue 740}
\addcontentsline{toc}{chapter}{Der Koeffizient 8pi und der Status von Issue 740}

{\small\noindent\textbf{Zusammenfassung.}\quad\ignorespaces
Dok.~374 hat die Identifikation $E_\star=E_P$ als Lesart \SM{} geführt und
daraus geschlossen, \#740 sei unter dieser Lesart geschlossen. Auf OPH-Seite
ist der Koeffizient inzwischen abgeleitet: aus zwei Beobachtern mit
verschiedener Spannung folgt $\kappa=b\,L^2/q$, mit dem modularen Koeffizienten
$b=2\pi$ und dem Viertel-Flächengesetz $q=1/4$ also $\kappa=8\pi L^2$ und
$G=L^2$ \QM. Damit ist $E_\star=\hbar c/L$ die nicht reduzierte Planck-Energie
\BM. Die weitergehende Aussage aus Dok.~374 wird zurückgenommen: OPH führt
Längenidentität, Uhr und einen Matching-Faktor weiterhin als offen, acht von
fünfzehn Abhängigkeiten von \#740 sind dort nicht geschlossen. Beide Rahmen
haben an derselben Stelle einen Matching-Faktor: OPH $0{,}16\,\%$ zwischen
Quell-$v$ und Fermi-$v$, FFGFT $1{,}10\,\%$ als bare-Rest der Leiter
(Dok.~375).
}\medskip

\section*{Zeichenerklärung}
\addcontentsline{toc}{section}{Zeichenerklärung}

\medskip
\begin{tabular}{@{}ll@{}}
\textbf{[B]} & algebraisch bewiesen \\
\textbf{[K]} & numerisch bestätigt \\
\textbf{[S]} & strukturell plausibel, Herleitung offen \\
\textbf{[Q]} & aus externer Quelle übernommen, zitiert \\
\end{tabular}

\medskip
\begin{tabular}{@{}lp{10cm}@{}}
\toprule
\textbf{Symbol} & \textbf{Bedeutung} \\
\midrule
$\kappa$ & Kopplung in der Einstein-Gleichung, $\kappa=8\pi G/c^4$ \\
$b$, $q$ & modularer Koeffizient und Vorfaktor des Flächengesetzes $S=A/(4L^2)$ \\
$L$ & Länge im Flächengesetz \\
$E_\star$ & OPH-Referenzenergie der Hierarchie \\
$E_P$, $\bar E_P$ & Planck-Energie, nicht reduziert und reduziert \\
$m$ & Matching-Faktor zwischen Quell- und Konventionsgröße \\
$\xi$ & FFGFT-Skalierungsparameter ($4/30000$) \\
\bottomrule
\end{tabular}

\section{Anlass}

Dok.~374 hat die Brückengleichungen zu OPH aufgestellt und die Identifikation
$E_\star=E_P$ als Lesart geführt. Im Ergebnisteil steht dort, \#740 sei unter
dieser Lesart geschlossen. Zwei Dinge haben sich seither geändert: Der
Koeffizient ist auf OPH-Seite abgeleitet, und OPH hat den Stand von \#740
ausführlicher dargestellt. Dieses Dokument hält beides fest und korrigiert die
zu weit gehende Aussage.

\section{Das Theorem \QM}

Zwei Beobachter lesen dieselbe Spannungs- und Geometriekonfiguration bei
gleicher Metriknorm. Die Tensorrelation lautet $F_i=\kappa\,T_i+\varphi\,n_i$
mit unbekanntem $\kappa$ und gemeinsamem Versatz $\varphi$. Auf der
Entropieseite stehen eine Volumenvariation $b\,K_i\,T_i$, eine
Flächenvariation $-K_i(F_i-\psi\,n_i)$ und die Stationarität
$\Delta S_i+(q/L^2)\,\Delta A_i=0$. Sehen die beiden Beobachter verschiedene
Spannung, $T_1\neq T_2$, so folgt ohne Annahme über die Versätze
\begin{equation}
  \kappa=\frac{b\,L^2}{q}.
\end{equation}
Mit $b=2\pi$ aus dem modularen Fluss der Kappe und $q=1/4$ aus dem
Flächengesetz ergibt das $\kappa=8\pi L^2$ \KM, und aus $\kappa=8\pi G$ folgt
$G=L^2$.

\section{Folge für die Referenzenergie}

Mit $L=\sqrt G$ (in Einheiten $\hbar=c=1$) ist
\begin{equation}
  E_\star=\frac{\hbar c}{L}=\sqrt{\frac{\hbar c^5}{G}}=E_P ,
\end{equation}
also die nicht reduzierte Planck-Energie \BM. Die reduzierte Lesart liegt um
$\sqrt{8\pi}=5{,}013$ daneben und verlangte eine andere Normierung des
Flächengesetzes; sie ist damit keine Einheitenwahl, sondern eine andere
physikalische Annahme.

Damit ist die Lesart aus Dok.~374 auf OPH-Seite nicht mehr gesetzt, sondern
abgeleitet. Der Vorfaktor $8\pi$ bleibt dort, wo er immer stand, nämlich in
der Einstein-Gleichung; das Viertel im Flächengesetz wählt $G$ aus.

\section{Korrektur zu Dok.~374}

Dok.~374 schließt daraus, \#740 sei unter der Lesart $E_\star=E_P$
geschlossen. Das ist zurückzunehmen. OPH führt \#740 als den gesamten
gemeinsamen Weltverbund und nennt acht von fünfzehn Abhängigkeiten als offen
\QM. Offen bleiben dort insbesondere:

\begin{enumerate}\itemsep2pt
\item Die Länge im Flächengesetz und die Länge im Pixel sind bislang durch
  Erklärung dieselbe; die Herleitung als eine Länge auf einer Quellfamilie ist
  geplant.
\item Der Takt. $\hbar c/\ell_P$ ist eine Energie; eine Uhr verlangt
  entwickelnden Zustand, Lücke und Auslesung sowie die Abbildung auf einen
  Laborübergang. Auf FFGFT-Seite liefert $\tilde T\cdot m=1$ den Takt direkt
  (Dok.~374, B3); auf OPH-Seite ist er in Arbeit.
\item Der Matching-Faktor $m=v_F/v_{\rm Quelle}=0{,}9984$, also die
  $0{,}16\,\%$. Er soll aus derselben Materiewirkung berechnet werden.
\end{enumerate}

Richtig ist damit: Der Koeffizient ist abgeleitet, die Brückengleichungen
B1--B4 aus Dok.~374 stehen unverändert, und die absolute Skala folgt über die
FFGFT-Kette. \#740 selbst bleibt offen.

\section{Matching-Faktoren beider Rahmen}

Beide Rahmen haben an derselben Stelle einen offenen Faktor, mit
verschiedener Herkunft:

\begin{center}
\begin{tabular}{@{}llrl@{}}
\toprule
\textbf{Rahmen} & \textbf{Faktor} & \textbf{Größe} & \textbf{Herkunft} \\
\midrule
OPH \QM & $v_F/v_{\rm Quelle}=0{,}9984$ & $0{,}16\,\%$ & Quell-$v$ gegen Fermi-$v$, Berechnung geplant \\
FFGFT & $v_F/v(\xi)=0{,}98912$ & $1{,}10\,\%$ & bare-Rest der Leptonleiter (Dok.~352, 375) \\
\bottomrule
\end{tabular}
\end{center}

Die FFGFT-Seite kennt die Herkunft ihres Faktors und führt ihn korpusweit als
bare gegen gemessen; die OPH-Seite führt ihren Faktor als offen. Für eine
Vergleichstabelle gehören beide Faktoren neben den Wert, wie der Rechenweg
auch (Dok.~375).

\section*{Werkzeug und Methode}
\addcontentsline{toc}{section}{Werkzeug und Methode}
Claude (Anthropic) wurde als Werkzeug eingesetzt für: Schreiben und Ausführen
des Python-Prüfskripts, Ausarbeiten der \LaTeX-Quellen und Pflege von Register
und Changelog. Das Werkzeug rechnet und schreibt; der Autor entscheidet, was
berechnet und geschrieben wird, beurteilt jedes Ergebnis und bestimmt, was es
physikalisch bedeutet. Das Theorem und die Angaben zum Stand von \#740 sind
aus der Mitteilung des OPH-Autors übernommen und als \QM{} gekennzeichnet.

\section{Ergebnis}

\begin{enumerate}
\item $\kappa=b\,L^2/q$ mit $b=2\pi$ und $q=1/4$ gibt $\kappa=8\pi L^2$ und
  $G=L^2$ \KM{} \QM.
\item Daraus ist $E_\star=\hbar c/L$ die nicht reduzierte Planck-Energie \BM;
  die reduzierte Lesart liegt um $\sqrt{8\pi}$ daneben.
\item Die Aussage aus Dok.~374, \#740 sei unter der Lesart geschlossen, wird
  zurückgenommen; Längenidentität, Takt und Matching-Faktor bleiben auf
  OPH-Seite offen \QM.
\item Matching-Faktoren: OPH $0{,}16\,\%$, FFGFT $1{,}10\,\%$ als bare-Rest
  \KM.
\end{enumerate}

\medskip
\noindent\textbf{Prüfskript:} \texttt{2/python/Dok376\_Skripte/pruef\_376\_koeffizient.py} (9/9)

\medskip
\noindent\textbf{Register (Dok.~190):} R124 (Rücknahme der \#740-Aussage in Dok.~374).

\begin{thebibliography}{9}
\bibitem{dok352} J.~Pascher, \emph{Dok.~352 --- Leptonreste} (2026).
\bibitem{dok374} J.~Pascher, \emph{Dok.~374 --- Gemeinsamer Anker} (2026).
\bibitem{dok375} J.~Pascher, \emph{Dok.~375 --- Die Hierarchie $v/E_P$} (2026).
\bibitem{oph} B.~Mueller et al., \emph{Observer Patch Holography},
  \texttt{github.com/FloatingPragma/observer-patch-holography}; Theorem und
  Statusangaben aus der Mitteilung vom 22.~September 2026.
\end{thebibliography}
